\documentclass[11pt,a4paper]{article}
\usepackage[T1]{fontenc}
\usepackage[utf8]{inputenc}
\usepackage[main=italian,provide=*]{babel}
\usepackage{amsmath,amssymb}
\usepackage{geometry}
\geometry{margin=2.3cm,top=2.5cm,bottom=2.7cm}
\usepackage{booktabs}
\usepackage[table,dvipsnames]{xcolor}
\usepackage{tcolorbox}
\tcbuselibrary{skins,breakable}
\usepackage{titlesec}
\usepackage{enumitem}
\usepackage{needspace}
\usepackage{fancyhdr}
\usepackage{graphicx}
\usepackage{hyperref}
\hypersetup{colorlinks=true,linkcolor=Sepia,citecolor=Sepia,urlcolor=Sepia}
\newcommand{\Tr}{\operatorname{Tr}}
\newcommand{\ket}[1]{\lvert #1 \rangle}
\newcommand{\bra}[1]{\langle #1 \rvert}

\definecolor{inkblue}{HTML}{1B3A5C}
\definecolor{lightblue}{HTML}{EAF1F8}
\definecolor{accent}{HTML}{B0562A}
\definecolor{softgray}{HTML}{6B6B6B}
\definecolor{okgreen}{HTML}{2E6B3E}
\definecolor{okgreenbg}{HTML}{EAF5EC}

\titleformat{\section}
  {\normalfont\Large\bfseries\color{inkblue}}
  {\thesection}{0.6em}{}[{\color{inkblue!40}\titlerule}]
\titlespacing{\section}{0pt}{0.9em}{0.4em}
\titleformat{\subsection}
  {\normalfont\large\bfseries\color{accent}}
  {}{0em}{}
\titlespacing{\subsection}{0pt}{0.6em}{0.2em}

\pagestyle{fancy}
\fancyhf{}
\renewcommand{\headrulewidth}{0.4pt}
\fancyhead[L]{\small\color{softgray}Scan sui parametri di rumore}
\fancyhead[R]{\small\color{softgray}\thepage}
\fancyfoot[C]{\small\color{softgray}Tesi triennale, Samuele --- Relatore: Prof.\ Alessandro Chiesa}

\newtcolorbox{keyeq}[1][]{
  colback=lightblue, colframe=inkblue!70, boxrule=0.6pt,
  arc=2pt, left=8pt, right=8pt, top=3pt, bottom=3pt, #1
}
\newtcolorbox{okbox}[1]{
  colback=okgreenbg, colframe=okgreen!60, boxrule=0.6pt,
  arc=2pt, left=7pt, right=7pt, top=2pt, bottom=2pt,
  fonttitle=\bfseries\color{okgreen}, title={#1}
}

\title{\vspace{-1.5em}\color{inkblue}Scan sui parametri di rumore\\[0.2em]
\Large\normalfont\color{softgray}Come si sposta $N^*$ al variare del canale}
\author{Note di lavoro --- Tesi triennale, Samuele\\ Relatore: Prof.\ Alessandro Chiesa}
\date{}

\begin{document}
\sloppy
\maketitle
\thispagestyle{fancy}
\vspace{-2.0em}

File di codice: \texttt{scan\_parametri\_rumore\_dimero.py}. Punto di
lavoro: $b/J=0.35$, $D/J=0.80$, $J=1$, $t=2$. Metrica: fedeltà rispetto
allo stato bersaglio fisico (quantum simulation Trotter sotto rumore).

\tableofcontents

\section{Obiettivo}

I passi precedenti hanno trovato $N^*$ a un \textbf{singolo} punto della
calibrazione di riferimento (\texttt{ibm\_torino}). Questo documento
chiede: $N^*$ si sposta se il rumore è più o meno forte? E se il
readout è più o meno rumoroso? Non è una domanda oziosa: un chip diverso,
o lo stesso chip in un momento diverso, avrà calibrazioni diverse ---
sapere come $N^*$ dipenda dai parametri del canale rende il risultato
utile oltre il singolo punto già misurato.

\section{Teoria: perché ci si aspetta un effetto su \texorpdfstring{$\varepsilon_{1q},\varepsilon_{2q}$}{eps} ma non su \texorpdfstring{$p_\text{readout}$}{p readout}}

Il rumore di gate entra nella fedeltà in modo non banale: più forte è
$\varepsilon_{1q}$ o $\varepsilon_{2q}$, più in fretta il rumore
accumulato supera il guadagno dell'errore di Trotter al crescere di $N$
--- ci si aspetta che $N^*$ \textbf{diminuisca} al crescere del rumore di
gate (meno passi convengono prima che il rumore domini).

Il readout, invece, entra come nella Sez.~3 del documento sui correlatori
dinamici: un fattore \textbf{puramente moltiplicativo},
$(1-2p_\text{readout})$, applicato a tutti gli $N$ allo stesso modo.
\begin{keyeq}
\centering
Se $g(N)$ ha un massimo a $N^*$, allora $c\cdot g(N)$ ($c>0$ costante)
ha il massimo \textbf{allo stesso} $N^*$: moltiplicare per una costante
positiva non sposta un argmax.
\end{keyeq}
Previsione, prima della verifica: $N^*$ dipende da $\varepsilon_{1q}$ e
$\varepsilon_{2q}$, non da $p_\text{readout}$.

\section{Risultato: monotonia rispetto a \texorpdfstring{$\varepsilon_{1q},\varepsilon_{2q}$}{eps}}
\label{sec:monotonia}

\begin{figure}[h]
\centering
\includegraphics[width=0.95\linewidth]{figure/fig_passo5_Nstar_vs_eps.pdf}
\caption{\textit{$N^*$ in funzione di $\varepsilon_{2q}$ (sinistra) e
$\varepsilon_{1q}$ (destra), scala log, l'altro parametro fissato al
valore di riferimento. Linea tratteggiata: valore di calibrazione
\texttt{ibm\_torino}. Monotonia non-crescente confermata in entrambi i
casi, con una transizione a $N^*=1$ per rumore sufficientemente forte
(il singolo passo di Trotter, per quanto impreciso, costa comunque meno
rumore accumulato di qualunque $N>1$).}}
\end{figure}

\begin{center}
\renewcommand{\arraystretch}{1.2}
\begin{tabular}{@{}cc@{}}
\toprule
$\boldsymbol{\varepsilon_{2q}}$ & $\boldsymbol{N^*}$\\
\midrule
$1.0\times10^{-3}$ & $10$\\
$2.0\times10^{-3}$ & $9$\\
$3.8\times10^{-3}$ (riferimento) & $8$\\
$6.0\times10^{-3}$ & $7$\\
$1.0\times10^{-2}$ & $1$\\
$1.5\times10^{-2}$ & $1$\\
$2.5\times10^{-2}$ & $1$\\
\bottomrule
\end{tabular}
\qquad
\begin{tabular}{@{}cc@{}}
\toprule
$\boldsymbol{\varepsilon_{1q}}$ & $\boldsymbol{N^*}$\\
\midrule
$1.0\times10^{-4}$ & $9$\\
$2.9\times10^{-4}$ (riferimento) & $8$\\
$1.0\times10^{-3}$ & $7$\\
$3.0\times10^{-3}$ & $1$\\
$1.0\times10^{-2}$ & $1$\\
$3.0\times10^{-2}$ & $1$\\
\bottomrule
\end{tabular}
\end{center}

\begin{okbox}{Verificato}
$N^*$ monotono non-crescente in entrambi $\varepsilon_{1q}$ e
$\varepsilon_{2q}$, come previsto. La transizione a $N^*=1$ non è un
artefatto numerico: oltre una certa soglia di rumore di gate, nessun
numero di passi di Trotter maggiore di 1 vale il costo aggiuntivo in
gate rumorosi.
\end{okbox}

\section{Risultato: invarianza rispetto a \texorpdfstring{$p_\text{readout}$}{p readout}}

\begin{center}
\renewcommand{\arraystretch}{1.2}
\begin{tabular}{@{}ccc@{}}
\toprule
$\boldsymbol{p_\text{readout}}$ & $\boldsymbol{N^*}$ & $\boldsymbol{F(N^*)}$\\
\midrule
$0.00$ & $5$ & $0.2621$\\
$0.05$ & $5$ & $0.2359$\\
$0.20$ & $5$ & $0.1572$\\
\bottomrule
\end{tabular}
\end{center}

\begin{okbox}{Verificato: la previsione teorica regge}
$N^*=5$ per tutti i valori di $p_\text{readout}$ testati, dallo $0\%$ al
$20\%$ --- un intervallo enormemente più ampio del valore di
calibrazione reale ($2.3\%$). Il valore assoluto $F(N^*)$ cambia
sensibilmente (da $0.26$ a $0.16$), ma la \emph{posizione} del massimo
no: esattamente la firma di un fattore moltiplicativo costante in $N$.
\end{okbox}

\section{Un affinamento: quanto è rapida la saturazione a \texorpdfstring{$N^*=1$}{N*=1}?}

La Sez.~\ref{sec:monotonia} mostra che $N^*$ decresce al crescere di $\varepsilon_{2q}$,
fino a saturare a $N^*=1$. Un modello ingenuo per l'errore totale,
$\varepsilon_\text{tot}(N)\sim a/N + g\,\varepsilon_{2q}N$ (Trotter
decrescente più rumore crescente, entrambi additivi), predice
l'ottimo per derivazione diretta:
\begin{keyeq}
\centering
$\dfrac{d\varepsilon_\text{tot}}{dN} = -\dfrac{a}{N^2} + g\varepsilon_{2q} = 0
\ \Longrightarrow\ N^* = \sqrt{\dfrac{a}{g\,\varepsilon_{2q}}} \ \propto\ \varepsilon_{2q}^{-1/2}$
\end{keyeq}
--- una pendenza $-0.5$ in scala log-log. Confrontando con i punti non
ancora saturati della Sez.~\ref{sec:monotonia} ($\varepsilon_{2q}=10^{-3}$ a
$6\times10^{-3}$, $N^*$ da $10$ a $7$):

\begin{center}
\renewcommand{\arraystretch}{1.2}
\begin{tabular}{@{}lc@{}}
\toprule
\textbf{Quantità} & \textbf{Valore}\\
\midrule
Pendenza log-log prevista ($\varepsilon_\text{tot}$ additivo) & $-0.5$\\
Pendenza log-log osservata (fit sui 4 punti non saturati) & $-0.195$\\
\bottomrule
\end{tabular}
\end{center}

\begin{okbox}{Non è una contraddizione: il modello additivo era per la
metrica sbagliata}
Lo scostamento fra $-0.5$ atteso e $-0.195$ osservato non indica un
errore --- indica che il modello $\varepsilon_\text{tot}(N)\sim a/N+g\varepsilon_{2q}N$
era pensato per un errore additivo su una singola quantità, mentre qui
l'osservabile è una \textbf{fedeltà}. Con rumore depolarizzante
accumulato su molti passi, una fedeltà si comporta più come un
decadimento (moltiplicativo, $F(N)\sim(1-\varepsilon_{2q})^{cN}$ per
qualche costante $c$) che come una somma di due termini indipendenti ---
la forma funzionale esatta di $N^*(\varepsilon_{2q})$ per una fedeltà
resta un affinamento possibile, non necessario per rispondere alla
domanda del relatore (``come si sposta $N^*$''), a cui la
Sez.~\ref{sec:monotonia}
risponde già direttamente con i dati.
\end{okbox}

\section{Cosa non è coperto da questo scan}

Lo scan varia un parametro alla volta, con gli altri fissati al valore
di riferimento --- non esplora l'intero spazio congiunto
$(\varepsilon_{1q},\varepsilon_{2q},p_\text{readout})$. Non copre inoltre
punti di lavoro fisici diversi da quello usato in tutta questa pipeline
($b/J=0.35$, $D/J=0.80$): se e quanto $N^*$ dipenda anche dal punto
fisico è una domanda verificata separatamente per l'estensione VQE
noise-aware (due punti di lavoro, \texttt{risultati\_vqe\_noise\_aware\_definitivo.tex}),
non qui.

\section{Considerazioni finali sulla pipeline di rumore}

Questo chiude la pipeline in cinque stadi: dal modello di rumore
(due canali, quattro numeri) fino a qui, ogni stadio ha aggiunto un
livello di struttura --- energia statica, poi dinamica (Trotter), poi un
osservabile a due tempi che richiede una misura reale (correlatori), poi
la sensibilità del risultato ai parametri stessi del rumore. Il filo
conduttore verificato ad ogni stadio: il readout simmetrico è sempre un
fattore moltiplicativo, mai additivo --- una proprietà che smette di
valere non appena il readout diventa asimmetrico (verificato
separatamente, \texttt{teoria\_readout\_asimmetrico.tex}), e che è alla
base di ogni argomento di invarianza usato in questa pipeline.

\begin{okbox}{Sintesi in una riga}
$N^*$ dipende in modo monotono e prevedibile dal rumore di gate, non
dipende affatto dal readout (simmetrico) --- entrambi i fatti verificati,
non solo previsti.
\end{okbox}

\needspace{4\baselineskip}
\section*{File prodotti}
\addcontentsline{toc}{section}{File prodotti}

\texttt{scan\_parametri\_rumore\_dimero.py}, questo documento.

\end{document}
